\chapter{約化態、糾纏熵與么正演化}

密度矩陣最初是描述不完全製備的自然方式：量子系統已經製備完成，但我們並不確切知道它是如何製備的。糾纏則賦予同一個數學物件更深一層的意義。即使複合系統的態已完全確知，子系統通常仍沒有自己的態向量。它的約化密度矩陣包含了所有可由局部測量預測的資訊。我們將以分量形式推導這個結論，並以兩個自旋態為例檢驗，再用熵區分乘積態與糾纏態。接著，我們將轉向第二個問題：量子態如何隨時間變化。內積守恆將引出么正性、厄米生成元，以及Schr\"odinger方程。

\section{密度矩陣與不完全製備}

假設一個電子被放在強磁場中，並讓它沿磁場方向排列。有人選定磁場方向並製備了這個電子，卻沒有告訴我們方向為何。因此，即使我們不知道自旋究竟處於哪個態，它仍可能處於完全確定的純態。我們可能知道各種可能製備方式的機率分布，也可能完全不知道。密度矩陣\(\rho\)正是描述這類不完全資訊的量子物件。

經典類比是相空間上的機率密度，或離散構型上的機率分布。量子物件則是一個具有三項基本性質的算符：
\begin{equation}
\rho^\dagger=\rho,
\qquad
\operatorname{Tr}\rho=1,
\qquad
\lambda_i\geq 0,
\end{equation}
其中\(\lambda_i\)是\(\rho\)的本徵值。厄米性使本徵值為實數，並確保存在一組正交歸一本徵基。正性使這些本徵值可作為機率，而歸一化條件給出
\begin{equation}
\sum_{i=1}^{n}\lambda_i=1.
\end{equation}
在\(\rho\)的本徵基中，\(\lambda_i\)是第\(i\)個本徵態的機率。

完全無知的情況，是希爾伯特空間中沒有任何方向受到偏好：
\begin{equation}
\rho_{\mathrm{max}}=\frac{1}{n}I.
\end{equation}
每個本徵值都是\(1/n\)。在另一個極端，完全確知的情況以歸一化態\(\ket{\psi}\)表示：
\begin{equation}
\rho_\psi=\ket{\psi}\bra{\psi}.
\end{equation}
這是投影到由\(\ket{\psi}\)張成的一維子空間的投影算符。

\begin{classroomqa}
\lecturetimestamp{00:06:47}
\audiencequestion{能否將某種製備方式描述為在兩個非正交態之間以五五開的機率選擇？例如，兩個自旋方向相差\(45^\circ\)？}
\lecturerresponse{這樣的製備方式確實可行，但這兩個權重通常不是所得密度矩陣的本徵值。要先將投影算符作加權求和，再將結果對角化。本徵值所具有的清楚機率意義，是相對於\(\rho\)的正交歸一本徵基而言，不一定適用於製備者所指定的那些非正交態。}
\end{classroomqa}

\begin{classroomqa}
\lecturetimestamp{00:07:47}
\audiencequestion{純態和混合態有何不同？}
\lecturerresponse{純態的密度矩陣只有一個非零本徵值，且該值等於\(1\)。混合態則有一個以上的非零本徵值。較早期的文獻使用德文詞\emph{Gemisch}；現代術語則直接稱為\emph{混合態}。}
\end{classroomqa}

我們來驗證純態投影算符的譜。由歸一化條件可得
\begin{equation}
\rho_\psi\ket{\psi}
=\ket{\psi}\braket{\psi}{\psi}
=\ket{\psi},
\end{equation}
因此，\(\ket{\psi}\)是本徵值為\(1\)的本徵向量。若\(\ket{\phi}\)與\(\ket{\psi}\)正交，則
\begin{equation}
\rho_\psi\ket{\phi}
=\ket{\psi}\braket{\psi}{\phi}
=0.
\end{equation}
所以，所有與\(\ket{\psi}\)正交的方向，其本徵值都是零。

\begin{figure}[htbp]
\centering
\includegraphics[width=0.72\textwidth]{lecture_08_figure_06.png}
\caption{純態密度矩陣的作用就像投影算符：\(\rho\ket{\psi}=\ket{\psi}\)。}
\lectureframe{00:09:32}
\end{figure}

\section{由跡計算期望值}

密度矩陣的操作規則是
\begin{equation}
\overline{M}=\operatorname{Tr}(\rho M),
\end{equation}
其中\(M\)是可觀測量。對兩個因子的乘積，跡的循環性質給出
\begin{equation}
\operatorname{Tr}(\rho M)=\operatorname{Tr}(M\rho),
\end{equation}
無論\(\rho\)與\(M\)是否對易，這個等式都成立。

在任意正交歸一基\(\{\ket{i}\}\)中，
\begin{equation}
\operatorname{Tr}X=\sum_i\bra{i}X\ket{i}.
\end{equation}
這個數值與基底的選擇無關。對純態而言，
\begin{align}
\overline{M}
&=\operatorname{Tr}\!\left(\ket{\psi}\bra{\psi}M\right) \\
&=\sum_i\braket{i}{\psi}\bra{\psi}M\ket{i} \\
&=\bra{\psi}M\left(\sum_i\ket{i}\bra{i}\right)\ket{\psi} \\
&=\bra{\psi}M\ket{\psi}.
\end{align}
利用恆等算符的完備性關係
\begin{equation}
\sum_i\ket{i}\bra{i}=I,
\end{equation}
便可將密度矩陣的規則還原為熟悉的純態期望值。

\begin{figure}[htbp]
\centering
\includegraphics[width=0.74\textwidth]{lecture_08_figure_07.png}
\caption{純態投影算符的跡計算，以及恆等算符的完備性關係。}
\lectureframe{00:14:25}
\end{figure}

\begin{classroomqa}
\lecturetimestamp{00:14:38}
\audiencequestion{左邊的量應該是\(\overline{M}\)，也就是\(M\)的平均值嗎？}
\lecturerresponse{是的。跡不是算符\(M\)本身，而是標量期望值\(\overline{M}\)。}
\end{classroomqa}

對一般的密度矩陣，選取其本徵基：
\begin{equation}
\rho\ket{j}=\lambda_j\ket{j}.
\end{equation}
在\(\rho\)與\(M\)之間插入一組完備態：
\begin{align}
\overline{M}
&=\sum_i\bra{i}\rho M\ket{i} \\
&=\sum_{i,j}\bra{i}\rho\ket{j}\bra{j}M\ket{i} \\
&=\sum_{i,j}\lambda_j\delta_{ij}\bra{j}M\ket{i} \\
&=\sum_j\lambda_j\bra{j}M\ket{j}.
\end{align}
這就是一般的加權平均：\(\lambda_j\)是第\(j\)個本徵態的機率，而\(\bra{j}M\ket{j}\)是態\(M\)中的期望值。

\begin{classroomqa}
\lecturetimestamp{00:18:12}
\audiencequestion{為什麼最後一個矩陣元的兩個標記都變成\(j\)？仍在求和的是哪個指標？}
\lecturerresponse{因為\(\rho\)在這組基底下是對角的，也就是\(\bra{i}\rho\ket{j}=\lambda_j\delta_{ij}\)，所以只有\(i=j\)的項會留下來。原本的雙重求和因此化為單一求和，可以稱為對\(j\)求和。虛設求和指標的名稱沒有物理意義。}
\end{classroomqa}

\section{熵與資訊的缺失}

對經典機率分布\(\{p_i\}\)，定義
\begin{equation}
S=-\sum_i p_i\log p_i,
\end{equation}
其中\(0\log 0\)取其極限值\(0\)。若某一機率等於\(1\)，其餘機率皆為零，則\(S=0\)：這個狀態已知。若\(n\)種可能性等機率，則
\begin{equation}
p_i=\frac{1}{n},
\qquad
S=-n\left(\frac{1}{n}\log\frac{1}{n}\right)=\log n.
\end{equation}
因此，當機率分布涵蓋更多重要的可能性時，熵會增加。粗略來說，使用自然對數時，\(e^S\)就是權重相近的有效狀態數。

\begin{figure}[htbp]
\centering
\includegraphics[width=0.68\textwidth]{lecture_08_figure_08.png}
\caption{機率分布的熵，以及等機率時的熵\(S=\log n\)。}
\lectureframe{00:21:08}
\end{figure}

對密度矩陣而言，將相同定義套用在其本徵值上：
\begin{equation}
S(\rho)=-\sum_i\lambda_i\log\lambda_i
=-\operatorname{Tr}(\rho\log\rho).
\end{equation}
眼前需要的是第一個表達式；跡的形式則凸顯結果與基底無關。

\section{複合系統與局部測量}

考慮兩個子系統\(A\)與\(B\)。基底向量記為\(\ket{a,b}\)。若\(A\)有\(N\)個基底態，而\(B\)有\(n\)個，複合系統便有\(Nn\)個基底態。兩個系統不必相同：一個可以是二態自旋，另一個則是三態原子。

這對子系統最一般的純態為
\begin{equation}
\ket{\Psi}
=\sum_{a,b}\psi(a,b)\ket{a,b},
\qquad
\sum_{a,b}|\psi(a,b)|^2=1.
\end{equation}
現在讓 Alice 測量只作用在子系統\(A\)上的可觀測量\(M_A\)。在複合希爾伯特空間中，對應的算符是\(M_A\otimes I_B\)。其矩陣元為
\begin{equation}
\bra{a',b'}M_A\otimes I_B\ket{a,b}
=\bra{a'}M_A\ket{a}\,\delta_{b'b}.
\end{equation}
\(B\)標記不受算符作用。

將期望值完整寫出：
\begin{align}
\expect{M_A}
&=\bra{\Psi}M_A\otimes I_B\ket{\Psi} \\
&=\sum_{a',b'}\sum_{a,b}
\psi^*(a',b')\psi(a,b)
\bra{a'}M_A\ket{a}\delta_{b'b} \\
&=\sum_{a',a,b}
\psi^*(a',b)\psi(a,b)
\bra{a'}M_A\ket{a}.
\end{align}

\begin{classroomqa}
\lecturetimestamp{00:29:12}
\audiencequestion{為什麼左矢上的標記有撇號，右矢上的卻沒有？}
\lecturerresponse{計算矩陣元之前，左矢與右矢各自獨立求和，因此需要使用不同的虛設標記。撇號用來區分左矢與右矢的指標。接著，算符與內積會決定哪些標記必須相同。}
\end{classroomqa}

固定\(a\)與\(a'\)，並對\(b\)求和。結果是在子系統\(A\)上的矩陣：
\begin{equation}
\boxed{
\rho_A(a,a')
=\sum_b\psi(a,b)\psi^*(a',b)
}.
\end{equation}
這就是子系統\(A\)的約化密度矩陣；等價地說，它是純態投影算符對子系統\(B\)取偏跡的結果：
\begin{equation}
\rho_A=\operatorname{Tr}_B\ket{\Psi}\bra{\Psi}.
\end{equation}
局部期望值因此寫成
\begin{equation}
\expect{M_A}
=\sum_{a,a'}\rho_A(a,a')\bra{a'}M_A\ket{a}
=\operatorname{Tr}_A(\rho_A M_A).
\end{equation}

\begin{figure}[htbp]
\centering
\includegraphics[width=0.76\textwidth]{lecture_08_figure_09.png}
\caption{對\(B\)的指標求和後，留下以兩個\(A\)指標表示的矩陣。}
\lectureframe{00:33:50}
\end{figure}

\begin{classroomqa}
\lecturetimestamp{00:35:39}
\audiencequestion{這個局域可觀測量在子系統 \(B\) 上是否就是恆等算符？}
\lecturerresponse{正是如此。\(M_A\otimes I_B\)可以改變 \(A\) 的座標，但不會改變 \(b\)，也不會以其他方式依賴它。因此，它在 \(B\) 空間中的矩陣元是 \(\delta_{b'b}\)。}
\end{classroomqa}

因此，只要實驗侷限於子系統\(A\)，所有預測都包含在\(\rho_A\)之中。Bob 不必測量自己的子系統。只要不使用\(B\)的自由度，Alice 的描述就必須採用密度矩陣形式。

\section{純整體可以有混合的部分}

整個系統處於純態\(\ket{\Psi}\)，但\(\rho_A\)一般會有多個非零本徵值。完全了解整體，不代表每個部分都有自己的態向量。貓與槍可以處於同一個糾纏態向量；此時，單獨的貓由約化密度矩陣描述。

\begin{classroomqa}
\lecturetimestamp{00:38:25}
\audiencequestion{\(\operatorname{Tr}(M\rho)\)可以換成\((\operatorname{Tr}M)(\operatorname{Tr}\rho)\)嗎？}
\lecturerresponse{不可以。乘積的跡通常不等於跡的乘積。循環性質允許\(\operatorname{Tr}(M\rho)=\operatorname{Tr}(\rho M)\)，但不能將跡因式分解。}
\end{classroomqa}

\begin{classroomqa}
\lecturetimestamp{00:38:49}
\audiencequestion{我們一開始只取基底態\(\ket{a,b}\)，還是取最一般的純態？}
\lecturerresponse{我們取的是複合系統最一般的純態：\(\ket{\Psi}=\sum_{a,b}\psi(a,b)\ket{a,b}\)。單一的\(\ket{a,b}\)只是一個基底向量。令人意外的是，即使這個複合系統的向量已完全確知，其約化態仍可能是混合態。}
\end{classroomqa}

這與經典力學形成鮮明對比。完全掌握複合系統的經典資訊，就能確定它的每個座標，也就能確定每個子系統。經典機率分布當然可以取邊際分布；但如果原始分布集中在相空間的一個點上，每個邊際分布也都會集中於一個確定值。量子純態則不一定如此。

\begin{classroomqa}
\lecturetimestamp{00:40:46}
\audiencequestion{對\(B\)求和，和對經典的位置--速度分布積分以消去速度，有何不同？}
\lecturerresponse{兩種操作看起來相似，但在純態極限下有所不同。完全已知的經典相空間點會同時確定位置與速度，因此任一邊際分布仍然確定。完全已知的糾纏量子態，卻能使任一子系統呈現非平凡的混合態。因此，偏跡顯示了「整體為純態便意味著各部分也為純態」這項經典推論在量子情況下不成立。}
\end{classroomqa}

\section{乘積態：特殊情況}

\(\rho_A\)何時會保持為純態？答案是波函數必須能因式分解：
\begin{equation}
\psi(a,b)=\phi(a)\chi(b).
\end{equation}
假設\(\phi\)與\(\chi\)各自已歸一化。將其代入約化密度矩陣可得
\begin{align}
\rho_A(a,a')
&=\sum_b\phi(a)\chi(b)\phi^*(a')\chi^*(b) \\
&=\phi(a)\phi^*(a')\sum_b|\chi(b)|^2 \\
&=\phi(a)\phi^*(a').
\end{align}
因此
\begin{equation}
\rho_A=\ket{\phi}\bra{\phi}.
\end{equation}
它有一個本徵值\(1\)，其餘本徵值皆為\(0\)，熵則為零。

\begin{figure}[htbp]
\centering
\includegraphics[width=0.74\textwidth]{lecture_08_figure_10.png}
\caption{對因式分解的波函數，對\(B\)求和就等於將\(\chi\)歸一化。}
\lectureframe{00:46:35}
\end{figure}

\begin{classroomqa}
\lecturetimestamp{00:44:17}
\audiencequestion{\(\phi(a)\)本身就是\(A\)系統的態嗎？}
\lecturerresponse{它是該態的波函數分量。更明確地說，\(\ket{\phi}=\sum_a\phi(a)\ket{a}\)；\(\ket{\chi}=\sum_b\chi(b)\ket{b}\)也有類似的表示。因式分解表示\(\ket{\Psi}=\ket{\phi}\otimes\ket{\chi}\)，也就是\(A\)與\(B\)之間沒有糾纏的乘積態。}
\end{classroomqa}

對每一個局部可觀測量，
\begin{equation}
\expect{M_A}
=\operatorname{Tr}(\rho_A M_A)
=\bra{\phi}M_A\ket{\phi}.
\end{equation}
子系統\(A\)的表現完全等同於純態\(\ket{\phi}\)。因此，對雙系統純態而言，\(\rho_A\)的熵可以檢驗乘積結構：熵為零表示乘積態，熵非零則表示存在糾纏。熵接近零代表糾纏較弱；接近最大值則代表在這個二分方式下，該態與任何乘積態都相距甚遠。

\section{單態：純整體與最大的局部無知}

考慮雙自旋單態
\begin{equation}
\ket{\Psi^-}
=\frac{1}{\sqrt{2}}
\left(
\ket{\uparrow\downarrow}
-\ket{\downarrow\uparrow}
\right).
\end{equation}
其非零波函數分量為
\begin{equation}
\psi(\uparrow,\downarrow)=\frac{1}{\sqrt{2}},
\qquad
\psi(\downarrow,\uparrow)=-\frac{1}{\sqrt{2}},
\end{equation}
而\(\psi(\uparrow,\uparrow)=\psi(\downarrow,\downarrow)=0\)。

對自旋\(2\)取偏跡。對角元素為
\begin{align}
\rho_1(\uparrow,\uparrow)
&=|\psi(\uparrow,\uparrow)|^2
 +|\psi(\uparrow,\downarrow)|^2
=\frac{1}{2}, \\
\rho_1(\downarrow,\downarrow)
&=|\psi(\downarrow,\uparrow)|^2
 +|\psi(\downarrow,\downarrow)|^2
=\frac{1}{2}.
\end{align}
非對角元素為
\begin{align}
\rho_1(\uparrow,\downarrow)
&=\psi(\uparrow,\uparrow)\psi^*(\downarrow,\uparrow)
 +\psi(\uparrow,\downarrow)\psi^*(\downarrow,\downarrow) \\
&=0,
\end{align}
其共軛也為零。因此
\begin{equation}
\rho_1
=
\begin{pmatrix}
\frac12&0\\[2pt]
0&\frac12
\end{pmatrix}
=\frac12 I.
\end{equation}

\begin{figure}[htbp]
\centering
\includegraphics[width=0.76\textwidth]{lecture_08_figure_11.png}
\caption{單態中一個自旋的約化密度矩陣之四個矩陣元素。}
\lectureframe{00:55:35}
\end{figure}

\begin{classroomqa}
\lecturetimestamp{00:55:58}
\audiencequestion{單一自旋的約化態，其熵是多少？}
\lecturerresponse{兩個本徵值都是\(1/2\)，因此
\[
S_1=-2\left(\frac12\log\frac12\right)=\log 2.
\]
這是二態子系統可能具有的最大熵。最大混合態的特徵是\emph{本徵值}相等，而不是在任意基底下矩陣元素相等。}
\end{classroomqa}

自旋沿任意軸的期望值都為零：
\begin{equation}
\expect{\bm{\sigma}\cdot\hat{\bm n}}
=\operatorname{Tr}\!\left(
\rho_1\,\bm{\sigma}\cdot\hat{\bm n}
\right)
=\frac12\operatorname{Tr}
\left(\bm{\sigma}\cdot\hat{\bm n}\right)
=0.
\end{equation}

\begin{classroomqa}
\lecturetimestamp{00:58:37}
\audiencequestion{為什麼\(\bm{\sigma}\cdot\hat{\bm n}\)對所有方向的跡都為零？}
\lecturerresponse{每個泡利矩陣的跡都為零。由於\(\bm{\sigma}\cdot\hat{\bm n}=n_1\sigma_1+n_2\sigma_2+n_3\sigma_3\)，利用跡的線性性，對任何單位向量\(\hat{\bm n}\)計算出的跡都為零。因此，沿任意選定軸測量自旋，得到平行或反平行結果的機率相同。}
\end{classroomqa}

整對自旋處於確定的純態，熵為零；單獨看其中任一個自旋，則是最大混合態。這不是因為製備紀錄不完整而造成的無知，而是糾纏所決定的局部態。

\section{等振幅乘積態}

現在考慮
\begin{equation}
\ket{\Phi}
=\frac12\left(
\ket{\uparrow\uparrow}
+\ket{\uparrow\downarrow}
+\ket{\downarrow\uparrow}
+\ket{\downarrow\downarrow}
\right).
\end{equation}
四個振幅都等於\(1/2\)，因此四項各自貢獻\(1/4\)，總和為一，態已歸一化。對自旋\(2\)取偏跡可得
\begin{equation}
\rho_1
=
\begin{pmatrix}
\frac12&\frac12\\[2pt]
\frac12&\frac12
\end{pmatrix}.
\end{equation}
其跡為\(1\)，但
\begin{equation}
\det\rho_1
=\frac14-\frac14
=0.
\end{equation}

\begin{classroomqa}
\lecturetimestamp{01:02:23}
\audiencequestion{這裡如何由跡與行列式求出兩個本徵值？}
\lecturerresponse{對\(2\times2\)矩陣而言，本徵值之和等於跡，乘積等於行列式。因此\(\lambda_1+\lambda_2=1\)，且\(\lambda_1\lambda_2=0\)，本徵值只能是\(\{1,0\}\)。約化態是純態，熵為零。}
\end{classroomqa}

\begin{classroomqa}
\lecturetimestamp{01:03:18}
\audiencequestion{這兩個沒有糾纏的自旋，實際上處於什麼狀態？}
\lecturerresponse{兩者都沿正\(x\)軸排列。從因式分解可以看出答案：
\[
\ket{\Phi}
=\frac{\ket{\uparrow}+\ket{\downarrow}}{\sqrt2}
\otimes
\frac{\ket{\uparrow}+\ket{\downarrow}}{\sqrt2},
\]
也可以透過計算泡利矩陣的期望值，以實驗方式確認。}
\end{classroomqa}

對\(z\)分量，
\begin{equation}
\sigma_3=
\begin{pmatrix}
1&0\\[2pt]
0&-1
\end{pmatrix},
\qquad
\operatorname{Tr}(\rho_1\sigma_3)=0.
\end{equation}
對\(x\)分量，
\begin{equation}
\sigma_1=
\begin{pmatrix}
0&1\\[2pt]
1&0
\end{pmatrix},
\end{equation}
以及
\begin{align}
\rho_1\sigma_1
&=
\begin{pmatrix}
\frac12&\frac12\\[2pt]
\frac12&\frac12
\end{pmatrix}
\begin{pmatrix}
0&1\\[2pt]
1&0
\end{pmatrix}
=
\begin{pmatrix}
\frac12&\frac12\\[2pt]
\frac12&\frac12
\end{pmatrix}, \\
\operatorname{Tr}(\rho_1\sigma_1)&=1.
\end{align}

\begin{figure}[htbp]
\centering
\includegraphics[width=0.70\textwidth]{lecture_08_figure_03.png}
\caption{等振幅約化態的\(\sigma_1\)矩陣計算。}
\lectureframe{01:05:13}
\end{figure}

因此\(\expect{\sigma_1}=1\)，而\(\expect{\sigma_2}=\expect{\sigma_3}=0\)。若振幅稍微偏離精確的因式分解形式，本徵值通常就會從\(\{1,0\}\)變成兩個非零值。由此產生的小熵，量度的是少量糾纏。

\section{這是誰的熵？}

剛才計算的是子系統的熵。以單態為例，
\begin{equation}
S(\rho_{12})=0,
\qquad
S(\rho_1)=S(\rho_2)=\log 2.
\end{equation}
整體是純態，每個部分卻是混合態。因此，糾纏熵不符合直觀的加法規則\(S_{12}=S_1+S_2\)。對雙系統純態而言，兩個約化密度矩陣的非零本徵值相同，所以兩者的熵相等，即使整體的熵為零。

\begin{classroomqa}
\lecturetimestamp{01:09:03}
\audiencequestion{每個密度矩陣，歸根究柢都是糾纏密度矩陣嗎？}
\lecturerresponse{在實際描述上，我們可以談一般的不確定性，而不追究缺失的紀錄由什麼承載。若將製備裝置、環境或觀察者納入系統，這種不確定性通常可以視為與未觀測自由度之間的糾纏。數學上，這個想法是嚴格成立的：每個密度矩陣都能在更大的希爾伯特空間中進行純化，但該純化無法唯一決定實際的製備歷程。}
\end{classroomqa}

\section{從離散的經典運動到量子運動}

現在我們暫時離開糾纏，轉而探討態如何隨時間變化。對於只有有限個構型的經典系統，可逆演化就是一種排列置換。每個狀態都有一個後繼狀態與一個前驅狀態；整個運動可分解為若干循環。

\begin{figure}[htbp]
\centering
\includegraphics[width=0.70\textwidth]{lecture_08_figure_12.png}
\caption{以離散構型之間的箭頭表示經典可逆演化。}
\lectureframe{01:11:40}
\end{figure}

\begin{classroomqa}
\lecturetimestamp{01:11:51}
\audiencequestion{若構型是離散的，為什麼時間步也必須是離散的？}
\lecturerresponse{嚴格的經典二態模型沒有中間構型可供連續路徑經過。若將一枚實體硬幣慢慢旋轉，就必須把它的朝向納入模型，作為額外的連續座標；這樣一來，它就不再是只有兩種構型的模型。量子二態系統則不同：即使正交歸一基底只有兩個向量，其希爾伯特空間仍包含連續多種疊加態。}
\end{classroomqa}

自旋\(\tfrac12\)態可以透過疊加，從\(\ket{\uparrow}\)連續地旋轉到\(\ket{\downarrow}\)。有限維不代表態向量只有有限多個。這使離散量子系統也能隨時間連續演化。

可逆的微觀動力學也會保留不同狀態之間的區別。兩個不同的經典初始條件不會合併成同一個微觀狀態，否則便無法辨別原本是哪個初始條件。哈密頓力學中的可逆性也伴隨著李維定理：相空間體積守恆。

\begin{classroomqa}
\lecturetimestamp{01:15:26}
\audiencequestion{那麼，在多個初始條件趨近同一穩態的非線性分岔中，情況又如何？}
\lecturerresponse{吸引子出現在有效的耗散描述中，因為這類描述已捨去環境自由度。基本的微觀哈密頓演化仍然可逆；資訊或許會轉移到未觀測的變數，但不同的完整狀態不會真的變成同一個完整狀態。}
\end{classroomqa}

量子力學的對應說法是么正性。正交的初態會演化成彼此正交的末態；更一般地說，所有內積都會保持不變。

\begin{classroomqa}
\lecturetimestamp{01:16:50}
\audiencequestion{兩個態必須受到完全相同的影響嗎？在沿\(x\)方向的磁場中，初始的\(z\)向上態和\(z\)向下態會如何變化？}
\lecturerresponse{同一個演化算符必須作用在兩個態上。在橫向磁場中，兩個自旋會進動，但它們的布洛赫向量始終方向相反，態向量也始終彼此正交。若施加不同的外部作用歷程，就等於在比較不同的演化算符。}
\end{classroomqa}

\section{線性與么正性}

量子時間演化的第一項原理是線性：
\begin{equation}
\ket{\psi(t)}=U(t)\ket{\psi(0)}.
\end{equation}
演化算符取決於經過的時間，並滿足
\begin{equation}
U(0)=I.
\end{equation}
量子力學本身具有線性，而這條規則也一樣，是由其實驗上的成功所支持。

\begin{figure}[htbp]
\centering
\includegraphics[width=0.70\textwidth]{lecture_08_figure_13.png}
\caption{時間\(t\)的態，是將線性算符\(U(t)\)作用於初態所得。}
\lectureframe{01:19:00}
\end{figure}

\begin{classroomqa}
\lecturetimestamp{01:19:01}
\audiencequestion{在\(t=0\)時，\(U(t)\)必須為何？}
\lecturerresponse{它必須是恆等算符。經過的時間為零時，不論初始向量為何，態都不會改變。}
\end{classroomqa}

第二項原理是內積守恆：
\begin{equation}
\braket{\phi(t)}{\psi(t)}
=\braket{\phi(0)}{\psi(0)}.
\end{equation}
當\(\phi=\psi\)時，這會保持態的歸一化，也就保持總機率。若初態彼此正交，演化後仍能完全區分。對任意態而言，它們之間的幾何關係也會完整保留。

左矢的演化方式為
\begin{equation}
\bra{\phi(t)}
=\bra{\phi(0)}U^\dagger(t),
\end{equation}
因此
\begin{align}
\braket{\phi(t)}{\psi(t)}
&=\bra{\phi(0)}U^\dagger(t)U(t)\ket{\psi(0)} \\
&=\braket{\phi(0)}{\psi(0)}.
\end{align}
由於這個等式對任意一對初始向量都成立，
\begin{equation}
\boxed{U^\dagger(t)U(t)=I}.
\end{equation}

\begin{figure}[htbp]
\centering
\includegraphics[width=0.74\textwidth]{lecture_08_figure_04.png}
\caption{右矢與左矢的演化，以及導出\(U^\dagger U=I\)的內積條件。}
\lectureframe{01:23:30}
\end{figure}

\begin{classroomqa}
\lecturetimestamp{01:22:10}
\audiencequestion{為什麼左矢會乘上\(U^\dagger\)，而不是\(U\)？}
\lecturerresponse{取厄米共軛會顛倒因子的順序，並將每個算符換成其厄米共軛：
\[
\bigl(U\ket{\phi}\bigr)^\dagger
=\bra{\phi}U^\dagger.
\]
這種對偶作用是保持內積一致所必需的。}
\end{classroomqa}

\begin{classroomqa}
\lecturetimestamp{01:24:34}
\audiencequestion{若\(U\)只是一個複數，\(U^\dagger U=1\)代表什麼？}
\lecturerresponse{這表示\(|U|=1\)，所以\(U=e^{i\theta}\)位於單位圓上。么正矩陣是純相位的算符類比：它保持長度與內積。}
\end{classroomqa}

\section{無窮小演化與哈密頓量}

令\(\epsilon\)為一個極短的時間間隔。由於\(U(0)=I\)，可寫出一階展開：
\begin{equation}
U(\epsilon)=I-i\epsilon h+O(\epsilon^2).
\end{equation}
負號與因子\(i\)是為了方便而採用的約定。取厄米共軛可得
\begin{equation}
U^\dagger(\epsilon)
=I+i\epsilon h^\dagger+O(\epsilon^2).
\end{equation}
么正性於是要求
\begin{align}
I
&=U^\dagger(\epsilon)U(\epsilon) \\
&=\left(I+i\epsilon h^\dagger\right)
\left(I-i\epsilon h\right)+O(\epsilon^2) \\
&=I+i\epsilon(h^\dagger-h)+O(\epsilon^2).
\end{align}
一階項必須為零：
\begin{equation}
\boxed{h^\dagger=h}.
\end{equation}
因此，無窮小生成元是厄米算符。

\begin{figure}[htbp]
\centering
\includegraphics[width=0.76\textwidth]{lecture_08_figure_14.png}
\caption{一階近似下的么正性條件要求無窮小生成元為厄米算符。}
\lectureframe{01:29:35}
\end{figure}

在上述暫用的約定中，\(h\)的單位是角頻率。恢復物理單位後，定義能量算符
\begin{equation}
H=\hbar h.
\end{equation}
這個厄米算符就是哈密頓量，其本徵值就是能階。

將無窮小演化算符作用在一個態上：
\begin{equation}
\ket{\psi(\epsilon)}
=\left(I-\frac{i\epsilon}{\hbar}H\right)\ket{\psi(0)}
+O(\epsilon^2).
\end{equation}
減去初態、除以\(\epsilon\)，再取極限：
\begin{equation}
\frac{\mathrm d}{\mathrm dt}\ket{\psi(t)}
=-\frac{i}{\hbar}H\ket{\psi(t)}.
\end{equation}
等價地，
\begin{equation}
\boxed{
i\hbar\frac{\mathrm d}{\mathrm dt}\ket{\psi(t)}
=H\ket{\psi(t)}
}.
\end{equation}
這就是一般形式的Schr\"odinger方程。

\begin{figure}[htbp]
\centering
\includegraphics[width=0.76\textwidth]{lecture_08_figure_15.png}
\caption{無窮小差商成為態對時間的導數。}
\lectureframe{01:32:18}
\end{figure}

哈密頓量取決於物理系統的設定。若外部磁場或其他控制條件改變，哈密頓量也可能隨時間改變。若系統設定不隨時間改變，\(H\)就是固定的厄米算符。

\begin{classroomqa}
\lecturetimestamp{01:34:30}
\audiencequestion{內積守恆是否表示，原本沒有糾纏的系統必須一直保持沒有糾纏？}
\lecturerresponse{不是。么正性保持的是完整態向量之間的內積，而不是特定劃分下的乘積結構。交互作用哈密頓量可以使乘積態演化為糾纏態。若么正變換無法做到這一點，測量就無法使系統與測量儀器產生關聯。}
\end{classroomqa}

\section{熵的單位}

\begin{classroomqa}
\lecturetimestamp{01:35:08}
\audiencequestion{糾纏有焦耳之類的物理單位嗎？}
\lecturerresponse{糾纏熵是資訊量度，不是能量。使用自然對數時，單位是奈特；使用以二為底的對數時，單位是位元。兩者都是無量綱的計數慣例。}
\end{classroomqa}

熱力學的記號可能會使這一點不易察覺，因為波茲曼常數用來換算以克耳文表示的溫度與能量：
\begin{equation}
\expect{K}_{\mathrm{monatomic}}
=\frac32 k_{\mathrm B}T.
\end{equation}
可逆熱量關係為
\begin{equation}
\delta Q_{\mathrm{rev}}=T\,\mathrm dS_{\mathrm{th}}.
\end{equation}
若資訊熵為
\begin{equation}
S_{\mathrm{info}}=-\sum_i p_i\log p_i,
\end{equation}
則傳統的熱力學熵為
\begin{equation}
S_{\mathrm{th}}=k_{\mathrm B}S_{\mathrm{info}}.
\end{equation}
令\(k_{\mathrm B}=1\)，就能以能量單位表示溫度，並使熵保持無量綱。如此一來，波茲曼常數便是換算因子，而非一種新的糾纏單位。歷史上將波茲曼的統計思想與愛因斯坦在1905年關於布朗運動的研究連結起來的補充說明，不會改變資訊理論的定義。

\section{能量本徵態與相位}

\begin{classroomqa}
\lecturetimestamp{01:39:54}
\audiencequestion{為什麼將厄米生成元\(H\)認定為能量？}
\lecturerresponse{么正性本身只能證明生成元是厄米算符。將生成元認定為能量，則是額外的物理主張：能量生成時間平移。對於不隨時間改變的哈密頓量，支配演化的同一算符也是守恆量，其本徵值就是定態能階。具體的哈密頓量則提供了這項認定的實驗內容。}
\end{classroomqa}

普朗克常數與頻率的關係為
\begin{equation}
E=h\nu=\hbar\omega,
\qquad
\omega=2\pi\nu,
\qquad
h=2\pi\hbar.
\end{equation}

\begin{classroomqa}
\lecturetimestamp{01:40:30}
\audiencequestion{演化方程中應該包含\(h\)、\(\hbar\)，還是還要多乘一個\(2\pi\)因子？}
\lecturerresponse{使用角頻率\(\omega\)時，方程為\(E=\hbar\omega\)，而Schr\"odinger方程中包含\(\hbar\)。使用一般頻率\(\nu\)時，等價關係為\(E=h\nu\)。\(2\pi\)因子完全可由\(\omega=2\pi\nu\)和\(h=2\pi\hbar\)這兩個關係解釋。}
\end{classroomqa}

假設初始態是能量本徵向量：
\begin{equation}
H\ket{\psi(0)}=E\ket{\psi(0)}.
\end{equation}

\begin{classroomqa}
\lecturetimestamp{01:43:18}
\audiencequestion{如果態仍是能量本徵向量，是否就表示它沒有改變？}
\lecturerresponse{它可能會乘上一個隨時間變化的複數因子，同時仍是具有相同本徵值的本徵向量。酉性會將該因子限制為相位。因此，即使右矢的相位持續旋轉，所代表的射線仍保持不變。}
\end{classroomqa}

寫成
\begin{equation}
\ket{\psi(t)}=f(t)\ket{\psi(0)}.
\end{equation}

\begin{figure}[H]
\centering
\includegraphics[width=0.70\textwidth]{lecture_08_figure_05.png}
\caption{能量本徵態至多乘上一個純量因子\(f(t)\)，仍沿著希爾伯特空間中的同一方向。}
\lectureframe{01:43:54}
\end{figure}

代入Schr\"odinger方程可得
\begin{align}
i\hbar\frac{\mathrm d f}{\mathrm dt}\ket{\psi(0)}
&=f(t)H\ket{\psi(0)} \\
&=E f(t)\ket{\psi(0)}.
\end{align}
約去兩邊共同的向量：
\begin{equation}
\frac{\mathrm d f}{\mathrm dt}
=-\frac{iE}{\hbar}f(t).
\end{equation}
在\(f(0)=1\)的條件下，
\begin{equation}
f(t)=e^{-iEt/\hbar},
\qquad
\ket{\psi(t)}=e^{-iEt/\hbar}\ket{\psi(0)}.
\end{equation}

\begin{figure}[H]
\centering
\includegraphics[width=0.76\textwidth]{lecture_08_figure_16.png}
\caption{能量本徵值方程，以及由此得到的相位演化。}
\lectureframe{01:46:10}
\end{figure}

\begin{classroomqa}
\lecturetimestamp{01:45:39}
\audiencequestion{相位\(e^{-iEt/\hbar}\)中，\(t\)的係數是什麼？}
\lecturerresponse{它就是角頻率，
\[
\omega=\frac{E}{\hbar}.
\]
等價地，\(E=\hbar\omega\)。若指數中的係數是實數，描述的會是增長或衰減率；此處指數中的虛數部分則描述酉演化下的相位旋轉。}
\end{classroomqa}

這個相位與頻率的關係，以及Schr\"odinger方程，是同一時間演化定律的兩種形式。這個方程具有一般性；自由粒子、處於勢能中的粒子，以及磁場中的自旋，各自的差異由其哈密頓量體現。下一步就是選定並求解這些哈密頓量。
